\begin{proposition}[Canonical bundles for all coprime multiplicities]\hypertarget{tag-01S5}{}\label{inv:B2.1}
\StatementRecord{inv:B2.1}{proved}{01S5}{\SeriesRef{inv:B1.1}, \SeriesRef{inv:B1.2}; \citet[Theorem 1.1 and Proposition 2.3]{ViaclovskyLog2026}}

Let $X=X_{3k,4l}(u)$ with admissible translation data as in
\citet[Theorem~1.1, p.~1, and Proposition~2.3, pp.~5--6]{ViaclovskyLog2026},
$\gcd(3k,4l)=1$, and any nonzero $u$ in its construction domain.
Then
\[
K_X=f^*\OO_B(-1)\otimes
\OO_X\bigl((k-1)S_1+(3l-1)S_2\bigr).
\]
Moreover $(h^{0,q})_{q=0}^3=(1,1,0,0)$, $h^{p,0}=0$ for $p>0$,
and $(h^{1,1},h^{1,2})=(2,1)$ for every such pair.
The relative comparison maps in Proposition~\ref{inv:B1.1}
identify the cohomological calculation for all these multiplicities.
\end{proposition}


\begin{proof}
\begin{enumerate}[label=\arabic*.,leftmargin=*]
\item The cited proposition supplies exactly the free cyclic quotients used
in \ref{inv:B1.1}; its Theorem 1.1 supplies the stated sphere topology for its
admissible choices. Proposition~\ref{inv:B1.2} computes the cohomology by the relative
cotangent comparison.
For the canonical bundle we check the natural evaluation
$f^*f_*\omega_{X/B}\to\omega_{X/B}$.
The invariant relative determinant has weights $a_1=2k$ and $a_2=l$
under the clockwise generators of orders $3k$ and $4l$. The pullback
of a canonical form is $F(s)\,ds\wedge e$, where the character of
$ds\wedge e$ is $\zeta_m^{1+a}$. Therefore its invariant coefficients
are $s^{m-1-a}$ times an invariant function and a nowhere vanishing
unit. Since the quotient of total spaces is \'etale, the evaluation
vanishes to order $m-1-a$ on the reduced fibre downstairs. These orders
are $k-1$ and $3l-1$. At the cusp the form
$d\log x_1\wedge d\log x_2\wedge dq$ is a nowhere vanishing canonical
frame in every smooth toric chart, so no cusp component occurs.

\item By \ref{inv:B1.1} and relative duality,
$f_*\omega_{X/B}=(R^2f_*\OO_X)^*=\OO(1)$.
The evaluation divisor just computed therefore gives
$\omega_{X/B}=f^*\OO(1)\otimes
\OO((k-1)S_1+(3l-1)S_2)$. Tensoring by $f^*K_B=f^*\OO(-2)$
proves the formula.
\end{enumerate}
\end{proof}